\chapter{Semiconductors, Diodes, and Transistors}
\label{ch:semiconductors}

\section{Doping and the p-n Junction}

Pure (\textbf{intrinsic}) silicon is a poor conductor. Deliberately
introducing trace impurities (\textbf{doping}) transforms its
conductivity in a controllable way: adding an element with one more
valence electron than silicon (phosphorus, arsenic) creates
\textbf{n-type} material with an excess of free electrons; adding an
element with one fewer valence electron (boron) creates \textbf{p-type}
material with an excess of ``holes'' (missing-electron sites that behave
as mobile positive charge carriers).

Where p-type and n-type materials meet, a \textbf{p-n junction} forms.
Electrons and holes diffuse across the junction and recombine near the
boundary, leaving a thin \textbf{depletion region} depleted of free
carriers and building a small internal electric field that opposes
further diffusion. This junction is the foundation of essentially every
semiconductor device in electronics.

\section{The Diode}

A \textbf{diode} is a single p-n junction packaged with two terminals,
the \textbf{anode} (p-side) and \textbf{cathode} (n-side). It conducts
current easily in one direction (\textbf{forward bias}, anode positive
relative to cathode) once the applied voltage exceeds a threshold ---
about 0.6--0.7\,V for silicon, 0.2--0.3\,V for germanium, and roughly
1.8--3.3\,V (varying by color/wavelength) for light-emitting diodes ---
and blocks current almost completely in the opposite direction
(\textbf{reverse bias}), up to a maximum reverse voltage rating beyond
which it breaks down.

\begin{center}
\begin{circuitikz}
\draw (0,0) to[battery1,l=$U$] (0,2) to[D,l=diode] (3,2) to[R,l=$R$] (3,0) -- (0,0);
\end{circuitikz}
\end{center}

\subsection{Diode Types and RF/Amateur Applications}

\begin{itemize}
\item \textbf{Rectifier diodes}: convert AC to pulsating DC in power
supplies (Chapter~\ref{ch:power-sources}).
\item \textbf{Zener diodes}: designed to operate safely in reverse
breakdown at a well-defined voltage, used as simple voltage references
or shunt regulators.
\item \textbf{Varactor (varicap) diodes}: reverse-biased junction
capacitance varies with applied voltage, used for electronic tuning in
VFOs and PLL synthesizers.
\item \textbf{Schottky diodes}: metal-semiconductor junction with a
lower forward voltage drop and much faster switching than a standard
silicon junction diode, widely used in RF detectors and mixers.
\item \textbf{PIN diodes}: behave as a variable resistor at RF when
forward-biased with varying DC current, used as RF switches and
attenuators (e.g., for T/R switching).
\item \textbf{Light-emitting diodes (LEDs)}: emit light when forward
biased, used for indicators and, in laser diode form, for optical
communication experiments.
\end{itemize}

\section{Bipolar Junction Transistors (BJTs)}

A \textbf{bipolar junction transistor} adds a second junction, forming
either an \textbf{NPN} or \textbf{PNP} sandwich of doped regions ---
\textbf{emitter}, \textbf{base}, and \textbf{collector}. A BJT is a
\textbf{current-controlled current source}: a small current flowing into
the base controls a much larger current flowing from collector to
emitter (NPN), with the ratio between them called \textbf{current gain},
$\beta$ (or $h_{FE}$):
\[
\boxed{I_C = \beta I_B}, \qquad I_E = I_B + I_C.
\]
Typical small-signal transistors have $\beta$ in the range of roughly
50--300. This current-controlled amplification is the basic mechanism
behind BJT amplifier and switching circuits.

\begin{center}
\begin{circuitikz}
\draw (0,0) node[npn](Q1){}
  (Q1.base) node[left]{B}
  (Q1.collector) node[above]{C}
  (Q1.emitter) node[below]{E};
\end{circuitikz}
\end{center}

\subsection{Common Amplifier Configurations}

\begin{itemize}
\item \textbf{Common-emitter}: input at base, output at collector,
emitter grounded (often through a small resistor). High voltage and
current gain, inverts the signal ($180\degree$ phase shift), moderate
input/output impedance --- the most common general-purpose amplifier
configuration.
\item \textbf{Common-base}: input at emitter, output at collector, base
grounded (at AC). Low input impedance, high output impedance, no phase
inversion, good high-frequency performance --- often used in RF
amplifier and mixer stages.
\item \textbf{Common-collector} (\textbf{emitter follower}): input at
base, output at emitter, collector grounded (at AC). Voltage gain near
unity, but high current gain and impedance transformation (high input
impedance, low output impedance) --- widely used as a buffer stage.
\end{itemize}

\section{Field-Effect Transistors (FETs)}

A \textbf{field-effect transistor} is a \textbf{voltage-controlled
current source}: a voltage applied to the \textbf{gate} controls current
flowing between \textbf{drain} and \textbf{source}, with (ideally) no
gate current at all --- giving FETs an extremely high input impedance,
in contrast to the current-driven BJT.

\begin{itemize}
\item \textbf{JFET} (junction FET): gate forms a reverse-biased junction
with the channel; normally ``on,'' with gate voltage used to pinch the
channel off.
\item \textbf{MOSFET} (metal-oxide-semiconductor FET): gate is insulated
from the channel by a thin oxide layer, giving even higher input
impedance than a JFET; available in \textbf{enhancement mode} (normally
off, turned on by gate voltage --- the dominant type in digital logic
and switching power supplies) and \textbf{depletion mode} (normally on).
\end{itemize}

MOSFETs, particularly power MOSFETs, are widely used in RF power
amplifiers and switching applications for amateur equipment because of
their high input impedance, ease of biasing, and good high-frequency
switching behavior.

\begin{safetynote}[Static sensitivity]
MOSFET gates are separated from the channel by an extremely thin
insulating layer that can be punctured by static discharge well below
the voltage a person can even feel. Handle MOSFETs and other
static-sensitive devices using proper ESD precautions (grounded
workstation, anti-static packaging) until they are soldered into a
circuit.
\end{safetynote}

\section{The Transistor as a Switch}

Beyond linear amplification, both BJTs and FETs are used as electronic
\textbf{switches}: driven fully ``on'' (\textbf{saturation} for a BJT,
low on-resistance for a MOSFET) or fully ``off'' (\textbf{cutoff}), a
transistor can control power to a load (a relay, a keying line, a
transmitter's PTT circuit) with a small control signal, forming the
building block of digital logic gates and switch-mode power circuits.

\section{Integrated Circuits}

An \textbf{integrated circuit (IC)} fabricates many transistors (from a
handful to billions), along with diodes, resistors, and sometimes small
capacitors, on a single piece of semiconductor material. This allows
complex functions --- operational amplifiers (Chapter~\ref{ch:opamps}),
digital logic, microcontrollers, RF mixers, phase-locked loop
synthesizers, and entire radio transceivers --- to be built far more
compactly, reliably, and cheaply than from discrete components, and
underlies essentially all modern amateur radio equipment design.

\begin{quickreview}
\item Doping creates n-type (excess electrons) and p-type (excess holes)
semiconductor material; a p-n junction is the basic building block.
\item Diodes conduct in one direction above a forward-voltage threshold
and block in reverse, up to a breakdown limit; many specialized diode
types (Zener, varactor, Schottky, PIN, LED) exploit particular junction
behaviors.
\item BJTs are current-controlled ($I_C = \beta I_B$); FETs are
voltage-controlled with very high input impedance.
\item Common-emitter, common-base, and common-collector configurations
trade off gain, impedance, and phase behavior differently.
\item Transistors are used both as linear amplifiers and as switches;
ICs integrate many transistors and other components on one chip.
\end{quickreview}

\begin{practiceproblems}
\item A BJT has $\beta = 100$ and $I_B = 50\,\mu\text{A}$. Find $I_C$ and
$I_E$.
\item Explain why a MOSFET's gate current is essentially zero in normal
operation, in terms of its construction.
\item Name one diode type suited to (a) voltage regulation, (b)
electronic tuning, (c) fast RF detection, and briefly justify each
choice.
\item Explain why an emitter-follower stage, despite having voltage gain
near 1, is still a useful amplifier configuration.
\item A technician reports a MOSFET failed immediately upon handling
before installation, with no circuit fault found. Suggest a likely
cause and a preventive measure.
\end{practiceproblems}
